\section{Boundaries}

V.6 does not say a voting system captured voter intent correctly. A CVR export
is an interpretation artifact. Hand audits, risk-limiting audits, recounts, and
adjudication records remain separate evidence.

V.6 also does not require every jurisdiction to publish ballot-level CVRs. Some
jurisdictions do not produce them, some publish only redacted or transformed
versions, and some contests are too small for safe ballot-level publication.
RCOUNT can still verify summary and manifest layers when public CVRs are absent.

The implemented row shape is a normalized interchange surface, not a direct
vendor schema. V.9 owns vendor CSV, JSON, XML, and NIST-style adapter details.
Adapters must preserve enough raw source evidence for an independent parser to
disagree with the normalized rows.

Finally, CVR publication must inherit V.4's receipt-safety rule. Rare ballot
styles, write-ins, timestamps, scanner batches, or unique selection patterns can
turn a CVR into identifying evidence. A deployable CVR layer needs redaction,
aggregation, or restricted-publication policy in addition to arithmetic checks.

\begin{center}
\small
\begin{tabularx}{\textwidth}{lX}
\toprule
CVR publication pattern & Risk \\
\midrule
Rare write-in plus tiny reporting unit & Narrows a ballot to a small group or one voter. \\
Unique cross-contest selection pattern & Can become a recognizable ballot signature. \\
Ballot style plus batch or scanner time & Can join with check-in or observation records. \\
Single-card provisional or cure batch & Can reveal whether a specific ballot was counted. \\
Joining CVRs to inclusion tokens & Can turn inclusion verification into a choice receipt. \\
\bottomrule
\end{tabularx}
\end{center}

Public CVR packages may therefore need aggregation, suppression, delayed
publication, or restricted evidence views. Hashes make evidence tamper-evident;
they do not make ballot-level records private.
